\chapter{Discusión}\label{cap:discusion}\label{cap:seleccion}
\section{Criterio de interpretación y selección}
Los resultados muestran que la selección de una estrategia de cuantización requiere evaluar conjuntamente el comportamiento de la aplicación y la implementación producida por las herramientas. En radiografías, varias variantes mantienen una exactitud próxima a la referencia, pero sus tiempos en placa difieren ampliamente. En MPC, las variantes INT8 reducen el tiempo de inferencia y la memoria, mientras que disminuyen de forma importante el cumplimiento del criterio de control. La configuración más rápida, por tanto, no coincide necesariamente con la alternativa que mejor conserva el comportamiento requerido.

La discusión se apoya en tres niveles de evidencia del Capítulo~\ref{cap:resultados}: desempeño del modelo en computador, concordancia y perfil de la implementación en placa, y funcionamiento del sistema integrado. La exactitud sobre 624 imágenes, la validación ST sobre ocho imágenes y la demostración UART con 16 imágenes corresponden a conjuntos diferentes. De igual forma, las 200 condiciones de simulación MPC, las 200 entradas de validación aislada y los ocho ensayos HIL responden a preguntas distintas. Sus porcentajes no se combinan como si procedieran de una única población experimental.

La selección se formula respecto de los objetivos efectivamente evaluados. Se favorece una implementación que conserve el comportamiento de interés y reduzca el costo de ejecución; cuando una mejora temporal implica una pérdida funcional, se explicita ese compromiso. No se introduce un umbral de aceptación clínica ni una garantía de estabilidad que la metodología no haya establecido. Tampoco se asigna una puntuación energética, pues no se dispone de mediciones de energía de estas ejecuciones.

\section{Radiografías: fidelidad de clasificación y costo de implementación}
\subsection{Qué se conserva y qué cambia con la cuantización}
La comparación sobre 624 imágenes muestra una exactitud de 79,81\% para Float32 y de 80,93\% para PTQ INT8. La diferencia de 1,12 puntos porcentuales favorece al archivo PTQ evaluado, pero no demuestra por sí sola que reducir la precisión mejore la capacidad de generalización. El contraste de grafos identificó una diferencia de topología: el PTQ INT8 desplegado presenta convoluciones de ocho canales y 30\,682 parámetros, mientras que la referencia y las variantes asociadas de seis canales presentan 22\,700 parámetros. Esta diferencia limita cualquier atribución causal exclusiva a la cuantización.

El análisis por clase también modifica la lectura de la exactitud. PTQ INT8 presenta una sensibilidad de 97,69\% y una especificidad de 52,99\%, frente a 98,21\% y 49,15\% en Float32. En términos de conteos, reduce los falsos positivos de 119 a 110, pero aumenta los falsos negativos de siete a nueve. La mejora de exactitud total no significa que todas las clases mejoren simultáneamente. La elección técnica debe conservar visible esta diferencia, especialmente porque el conjunto contiene más imágenes de neumonía que normales.

Mixed-5 conserva exactamente las 624 decisiones de Float32 y presenta un RMSE de salida de 0,003778. Ese resultado respalda su fidelidad respecto de la referencia en el conjunto observado, aunque no demuestra identidad de las puntuaciones ni equivalencia sobre imágenes externas. Su interés experimental reside en mantener las decisiones con una representación parcialmente cuantizada. QAT obtiene 79,97\% y el menor archivo TFLite del grupo. La validación en placa completa la comparación: informa 14,814 ms en M55 y 0,290 ms en NPU. La mejora temporal de NPU se acompaña de un falso positivo adicional en el lote validado, por lo que la exactitud del archivo en computador no se asigna automáticamente a esa ruta de destino.

\subsection{Latencia, memoria y asignación CPU/NPU}
Para el mismo PTQ INT8, la ruta Neural-ART reduce la media a 0,514 ms, una razón de tiempos de 32,34 respecto de M55. La concordancia cruzada presenta RMSE 0,002584 en las ocho entradas validadas. Esta combinación de baja latencia y diferencias de salida pequeñas respalda el aprovechamiento del acelerador en esa implementación concreta. El perfil atribuye 66,6\% del tiempo a hardware, 31,0\% a software y 2,3\% a control software: el beneficio global no exige que todo el programa se ejecute dentro de la NPU.

Mixed-5 ilustra un compromiso distinto. Su ruta M55 tarda 52,481 ms y la ruta Neural-ART, 104,700 ms; ambas presentan RMSE cruzado informado como cero en su validación. El tiempo de M55 es aproximadamente la mitad. En Neural-ART, el perfil asigna 99,6\% a software y las dos convoluciones software acumulan 92,836 ms, equivalentes al 88,67\% del total. La preservación de regiones flotantes conserva fidelidad, pero restringe el beneficio práctico del acelerador en este grafo.

La memoria refuerza el interés de M55 para Mixed-5: informa 26\,000 B de pesos y 533\,652 B de activaciones, frente a 26\,289 B y 781\,536 B en Neural-ART. En esta comparación, la ruta M55 combina menor tiempo y menores recursos de red con la concordancia numérica observada. Para PTQ INT8, en cambio, la aceleración NPU se acompaña de un aumento de activaciones de 92\,496 a 112\,500 B. La elección del motor debe considerar ambas magnitudes, sin confundirlas con el consumo total de Flash y SRAM del firmware.

\subsection{Justificación de la alternativa final para radiografías}\label{sec:decision_rx}
Bajo la preferencia por una menor discrepancia cruzada respecto de QAT/NPU, se propone \textbf{PTQ INT8 sobre Neural-ART} como candidato de despliegue provisional. Esta preferencia no establece un óptimo ni una superioridad general entre estrategias. La variante QAT/NPU es más rápida, pero presenta un RMSE cruzado mayor y una decisión adicional incorrecta en el lote de ocho imágenes. La propuesta se fundamenta en sus 0,514 ms de inferencia y la concordancia cruzada observada. El desempeño medido en PC caracteriza el archivo allí evaluado y no se traslada automáticamente al artefacto de placa sin acreditar su correspondencia. Es una selección técnica dentro del estudio, condicionada a la diferencia de topología y a que la validación de placa sólo comprende ocho entradas; no constituye una validación clínica ni acredita en placa la exactitud del conjunto completo de 624 imágenes.

Si se prioriza exclusivamente la menor latencia observada, \textbf{QAT sobre Neural-ART} pasa a ser el candidato temporal, con 0,290 ms y 98,9\% del tiempo en hardware. Su selección definitiva requiere resolver o aceptar explícitamente la diferencia numérica observada; no basta el perfil favorable. En ejecución sólo CPU, QAT/M55 ofrece 14,814 ms y RMSE cruzado informado como cero, además de menores pesos y activaciones que PTQ INT8/M55. Por ello constituye una alternativa pertinente cuando se busca concordancia con el archivo QAT y ejecución íntegra en CPU. Estas conclusiones conservan la limitación de comparar topologías distintas entre QAT y PTQ INT8.

\textbf{Mixed-5} conserva las decisiones Float32 en las 624 imágenes evaluadas en PC. Su implementación en \textbf{Cortex-M55} presenta concordancia con el archivo importado en el lote de placa validado y ofrece menor latencia y memoria que su ruta Neural-ART. No se ha acreditado en placa la conservación de las 624 decisiones; esa comprobación permanece pendiente. La demostración UART realizada con Mixed-5 conserva su valor como comprobación de la integración y de la comunicación. Sus 16 imágenes no convierten, sin embargo, a la configuración demostrada en la más rápida de las alternativas. Justificar el prototipo construido y recomendar el mejor despliegue según un criterio son decisiones relacionadas, pero diferentes.

La ruta radiográfica importada como INT16x8 en Neural-ART no se selecciona pese a sus 0,600 ms. El RMSE cruzado de 0,317093 y la concentración de salidas alrededor de 0,5 impiden considerar ese tiempo una aceleración numéricamente equivalente. Además, la igualdad de exactitud en ocho entradas oculta matrices de confusión diferentes. Este caso muestra por qué la validación de salidas debe preceder a la clasificación de alternativas por velocidad.

\FloatBarrier
\section{MPC: desempeño de control y elección del motor de ejecución}
\subsection{La fidelidad local no sustituye el comportamiento en lazo cerrado}
El controlador Keras obtiene 199 éxitos de 200 condiciones iniciales y la variante INT16x8, 198. PTQ INT8 por tensor, PTQ por canal y QAT alcanzan 74, 50 y 65, respectivamente. La precisión mixta recupera 156 éxitos. Estos resultados señalan que este caso es más sensible a las modificaciones numéricas que el clasificador, bajo las métricas y bancos de prueba utilizados. No permiten extrapolar una sensibilidad universal a toda red de control o a toda CNN.

Las curvas de posición, ángulo y acción de control ayudan a interpretar esa diferencia. Una posición acotada no basta para concluir recuperación del péndulo; el ángulo puede abandonar la banda final o experimentar giros amplios. Asimismo, una trayectoria particular que cumple el criterio no representa una tasa elevada sobre las 200 condiciones. La incorporación de los gráficos angulares permite relacionar el resultado agregado con los transitorios y con la pérdida tardía de recuperación observada en algunos ejemplos.

La validación aislada distingue dos errores que no deben confundirse. El RMSE cruzado compara la implementación de placa con su archivo de referencia; no mide directamente la diferencia entre ese archivo cuantizado y la política Float32, ni demuestra estabilidad del lazo. Por ello, una implementación INT8 de M55 puede presentar un RMSE cruzado informado como cero y, a la vez, obtener una tasa de éxito baja en simulación. En ese caso la placa reproduce el modelo evaluado, pero reproducirlo fielmente no corrige el comportamiento introducido al cuantizarlo.

La transferencia EQX--Keras también mostró que acciones casi idénticas ante entradas comunes no garantizan trayectorias idénticas al cerrar el lazo. Esta observación es consistente con la necesidad, expuesta en el marco teórico, de evaluar la evolución completa y las restricciones. Salamati et al. relacionan precisión de implementación y márgenes de robustez en MPC explícito \cite{salamati2019}; su garantía corresponde a otro controlador y procedimiento. Aquí se aporta evidencia experimental, sin trasladar esa garantía formal a la política neuronal.

\subsection{Compromiso entre precisión, memoria y tiempo}
La variante mixta M55 combina cálculo entero vectorial y fragmentos densos Float32 escalares. Su media de 0,326 ms supera en 7,59\% los 0,303 ms de Float32, aunque sus pesos se reducen de 62\,404 a 47\,404 B, un 24,04\%. Las activaciones aumentan de 400 a 852 B. Este resultado evidencia un compromiso entre almacenamiento de pesos, memoria de trabajo y tiempo; la reducción del ancho de ciertos tensores no garantiza una mejora simultánea de las tres magnitudes.

Los listados Float32 e INT16x8 para M55 no contienen las familias MVE examinadas ni operandos Q, y los informes de generación consignan el mismo volumen de pesos y activaciones. Junto con su latencia de 0,303 ms, esta evidencia respalda que la ruta obtenida desde INT16x8 no constituye una implementación nativa del esquema solicitado. La alta tasa de éxito del archivo INT16x8 en simulación sigue siendo relevante, pero no acredita un beneficio de cálculo INT16x8 en la placa. La capacidad de Helium y el soporte de una representación en la cadena de herramientas son propiedades distintas.

\subsection{Por qué Neural-ART no es la elección preferente para el modelo mixto}
Las variantes enteras alcanzan 0,067 ms en Neural-ART y un perfil con 98,7\% del tiempo en hardware. La aceleración es efectiva respecto de sus rutas M55, aunque la rapidez no compensa las tasas de éxito de 25--37\% del banco de simulación. La fidelidad funcional debe resolverse antes de preferir una variante por su latencia.

Para precisión mixta, Neural-ART requiere 0,480 ms frente a 0,326 ms en M55: un aumento de 47,24\%. El perfil distribuye 9,5\% en hardware, 84,1\% en software y 6,4\% en control software. Los grafos y los extractos del código sitúan regiones flotantes, conversiones y ejecución software dentro de la ruta generada. Esa distribución muestra que la ejecución medida está dominada por software y es consistente con la ausencia de ventaja temporal de esta ruta para el modelo seleccionado.

El fragmento Float32 vectorial encontrado en la ruta Neural-ART no contradice esa conclusión. La existencia de un núcleo MVE no obliga a que el programa completo sea más rápido que otra implementación con partición diferente. En esta comparación, la ruta M55 tiene menor latencia y menor RMSE cruzado informado. Neural-ART reduce modestamente los pesos, de 47\,404 a 47\,201 B, y las activaciones, de 852 a 616 B. Existe, por tanto, un intercambio entre tiempo y memoria; no una superioridad absoluta en todas las métricas.

La documentación de ST condiciona el soporte hardware al formato y a los operadores admitidos, y asigna operaciones flotantes a software \cite{cap4_stoperators}. Esta restricción resulta coherente con la fracción software medida, pero la conclusión sobre cuál ruta conviene procede de los perfiles de esta implementación. El nombre de un archivo cuantizado o la presencia de la opción NPU no sustituyen esa comprobación.

\subsection{Justificación de la alternativa final y alcance del HIL}\label{sec:decision_mpc}
Para continuar el objetivo experimental de desplegar una política \emph{parcialmente cuantizada}, se selecciona \textbf{precisión mixta sobre Cortex-M55} como compromiso entre recuperación funcional, reducción de pesos y costo temporal. Su 78\% de éxito supera a las variantes INT8 evaluadas; conserva regiones INT8 identificables y evita el incremento temporal de la ruta Neural-ART del mismo modelo. Esta elección justifica la alternativa investigada en el banco HIL, pero no demuestra que sea la mejor política de control de todas las comparadas.

Si el criterio principal fuera preservar el mayor desempeño de la referencia sin exigir cuantización, \textbf{Float32 sobre Cortex-M55} seguiría siendo la opción de referencia preferente: combina 99,5\% de éxito en simulación y 0,303 ms en la validación aislada. Los registros no contienen una campaña HIL equivalente para confirmar esa comparación en el sistema integrado. INT16x8 mantiene 99\% en simulación, pero la ruta compilada observada no aporta una ejecución nativa que justifique escogerla por una supuesta aceleración de enteros de 16 bits.

Los ocho ensayos HIL de la alternativa mixta obtuvieron cinco cumplimientos y tres incumplimientos del criterio final. El 62,5\% describe este conjunto determinista y no es una estimación directamente comparable con el 78\% del banco de 200 condiciones. Todos los ensayos respetaron los límites globales de posición y acción de control registrados, pero eso no implica cumplir la recuperación final del ángulo y la posición. En particular, recuperar el equilibrio después de un giro amplio tampoco equivale a mantenerlo durante todo el transitorio.

Las mediciones de inferencia y comunicación respaldan la viabilidad temporal del intercambio evaluado: el máximo RTT fue de 5,3675 ms frente a un período simulado de 160 ms. Esa holgura no demuestra ejecución estricta en tiempo real de todo el banco: siete registros de consola informan ciclos tardíos del computador y el octavo no permite concluir ausencia de retrasos. Tampoco se identifica de manera inequívoca el binario activo en cada CSV mediante un hash. La asociación del ensayo a la alternativa seleccionada procede del procedimiento experimental y sus registros; los tiempos compatibles no sustituyen esa identificación.

En consecuencia, se mantiene la precisión mixta M55 como \emph{alternativa experimental seleccionada}, con recuperación parcial demostrada y fallos explícitos. No se la presenta como un controlador que satisface todas las condiciones ensayadas. Si el requisito fuese cumplir el criterio en todo el banco, la evidencia disponible no permitiría cerrar esa aceptación para el modelo mixto; sería necesario mejorar su desempeño y repetir la evaluación antes de formular una conclusión de robustez general.

\FloatBarrier
\section{Qué aporta la comparación con NUCLEO-L476RG}\label{sec:discusion_l476}
\subsection{Comparación de plataformas y alcance de la evidencia}
La N657 presenta menor latencia en todos los modelos comunes: el cociente L476/M55 es 19,53 en radiografías QAT y entre 4,99 y 9,37 en MPC. El clasificador requiere aproximadamente la mitad de ciclos en M55, mientras que los cinco modelos MPC requieren más ciclos y terminan antes por la frecuencia diez veces mayor. La comparación caracteriza implementaciones completas; no aísla el efecto de Helium frente a arquitectura, memoria y bibliotecas.

Los listados aportan una explicación complementaria: documentan MVE en los núcleos M55 identificados y DSP en cinco de los seis ejecutables M4. El uso de DSP en inferencia está documentado en CMSIS-NN \cite{lai2018}; el estudio controlado de Arm respalda el potencial de MVE \cite{res_helium2021}. Estos antecedentes no permiten trasladar un factor de aceleración a nuestras mediciones. La evidencia detallada se conserva en el Anexo~\ref{anexo:instrucciones}; la selección se apoya en el desempeño funcional y los costos medidos.

\subsection{La memoria también determina qué alternativa puede compararse}
El clasificador QAT completa la validación en L476, mientras que los intentos PTQ INT8 y mixto fallan en etapas diferentes por restricciones de memoria de la configuración. Ello convierte a QAT en la única alternativa radiográfica validada para esa placa dentro de esta campaña. Su selección en L476 se justifica por factibilidad observada y por concordancia numérica con la referencia en las ocho imágenes; sus 289,351 ms no deben compararse con una latencia inexistente de los intentos fallidos.

Los gráficos de memoria muestran por qué el tamaño de los pesos no basta para evaluar despliegue. Las activaciones del clasificador tienen un peso importante, mientras que las variantes MPC necesitan espacios de trabajo mucho menores. Además, la implementación puede modificar esos búferes aunque los pesos y el hash de origen coincidan. La asignación física, el código, la pila y los búferes de validación también deben caber en el proyecto; la suma parcial de pesos y activaciones no certifica esa condición.

La ampliación no invalida la selección mixta M55 para el experimento MPC: sigue siendo más rápida que la ruta mixta NPU y conserva la ventaja funcional previamente observada frente a las variantes INT8. Sin embargo, L476 también puede ejecutar esa inferencia en 2,401 ms. Ese valor permite considerar una plataforma más pequeña cuando el presupuesto temporal y la memoria lo permitan, pero no acredita todavía un control HIL equivalente, una ventaja energética o una decisión de costo, dimensiones que no se midieron comparativamente.

\FloatBarrier
\section{Síntesis de las decisiones y respuesta al objetivo de la memoria}
\begin{table}[H]\centering\small
\caption{Selección justificada según el objetivo de cada aplicación.}\label{tab:decision_final}
\begin{tabularx}{\linewidth}{p{2.7cm}p{3.2cm}Y}\toprule
Objetivo & Alternativa priorizada & Fundamento y alcance\\\midrule
Radiografías: compromiso temporal y numérico & PTQ INT8 / Neural-ART & 0,514 ms y menor RMSE cruzado que QAT/NPU. Comparación condicionada por topología y lote reducido.\\
Radiografías: mínima latencia medida & QAT / Neural-ART & 0,290 ms; 5/8 aciertos frente a 6/8 de referencia. Requiere evaluar la diferencia de decisiones antes de adoptarlo.\\
Radiografías: conservar decisiones Float32 en PC & Mixed-5; candidata M55 & Coincidencia en 624 imágenes en PC; validación de placa en lote reducido. 52,481 ms y menor memoria que su ruta Neural-ART.\\
MPC: alternativa parcialmente cuantizada & Mixta / M55 & 78\% en simulación, 0,326 ms y reducción de pesos; cinco de ocho criterios HIL cumplidos. Selección experimental con fallos identificados.\\
MPC: priorizar desempeño de referencia & Float32 / M55 & 99,5\% en simulación y 0,303 ms. Falta una campaña HIL comparable para cerrar la evaluación integrada.\\\bottomrule
\end{tabularx}
\fuente{interpretación conjunta de las tablas y los perfiles del Capítulo~\ref{cap:resultados}}
\end{table}

La comparación con L476 y los dos casos en N657 respaldan una conclusión común: conviene seleccionar la combinación de modelo, precisión y ruta de ejecución. La NPU aporta una ventaja clara para el clasificador INT8 y reduce la latencia de las variantes enteras del controlador. Cortex-M55 constituye una alternativa pertinente para los modelos mixtos estudiados, porque evita una ruta Neural-ART dominada por software y dispone de núcleos vectoriales identificados en los listados. La ventaja concreta de M55 se acredita con los tiempos medidos; los listados acreditan la presencia estática de instrucciones Helium en funciones asociadas a la red, sin cuantificar su ejecución ni su contribución temporal durante los ensayos. Ambos tipos de evidencia se complementan, sin convertirlos en una medición aislada del efecto de MVE.

La precisión mixta ofrece una vía para conservar regiones sensibles y mantener parte de la red cuantizada, pero no es una solución universal. En radiografías puede conservar decisiones sin ofrecer la menor latencia. En MPC mejora el resultado frente a INT8, aunque no recupera el nivel Float32 y puede aumentar el tiempo de inferencia. El objetivo de evaluar estrategias y seleccionar por aplicación se resuelve, por tanto, mediante decisiones condicionadas a requisitos explícitos y al comportamiento observado, no mediante una preferencia general por un formato o acelerador.

La evidencia también delimita las expectativas iniciales del estudio: un modelo denominado INT16x8 no garantiza cálculo INT16x8 en el destino; un programa NPU puede ejecutar gran parte de su tiempo en software; y la coincidencia numérica con el archivo cuantizado no garantiza un desempeño adecuado en lazo cerrado. La contribución experimental consiste en documentar esas diferencias a lo largo del proceso completo, desde el archivo exportado hasta la aplicación integrada.

\section{Limitaciones y continuidad experimental}
Las conclusiones se circunscriben a las redes, versiones de herramientas y configuraciones evaluadas. En radiografías, la diferencia de topología, el tamaño reducido de las validaciones de placa y la cobertura incompleta de todas las variantes y placas impiden establecer una jerarquía universal. No se realizaron pruebas clínicas ni una validación externa del clasificador. En MPC, los ensayos HIL utilizan una planta simulada y no incluyen todos los modelos cuantizados ni todas las perturbaciones de una planta física.

La auditoría de 23 listados aumenta la trazabilidad de Helium, pero sigue siendo estática. No mide frecuencia de ejecución, ocupación del procesador o ahorro energético, ni sustituye el perfil por operador. No se dispone del listado radiográfico INT16x8 M55. Las diferencias pequeñas de tiempo tampoco permiten concluir significación estadística, pues no se evaluaron compilaciones repetidas.

Como trabajo futuro se propone comparar las cuantizaciones de radiografías sobre una misma topología y extender la validación de placa al conjunto común; en MPC, contrastar Float32 y precisión mixta mediante las mismas condiciones HIL y revisar los casos que incumplen el criterio. Una comparación del mismo ejecutable funcional con y sin MVE permitiría aislar el aporte de Helium. La medición energética y la identificación del firmware en cada corrida completarían dimensiones que actualmente permanecen abiertas. Estas acciones quedan fuera de la evaluación realizada.

La discrepancia de decisiones QAT/NPU no tiene una causa aislada demostrada en esta campaña; motiva ampliar la validación, sin establecer un umbral de rechazo no definido. Del mismo modo, las diferencias entre condiciones HIL de signo opuesto son observaciones de esos ensayos y no prueban un defecto de normalización. Los umbrales y las regiones de precisión mixta corresponden a las configuraciones ensayadas: no se realizó una búsqueda exhaustiva ni un análisis de sensibilidad que permita calificarlas como óptimas.

\section{Cumplimiento de los objetivos y conclusión general}\label{sec:cierre_objetivos}
El objetivo general se aborda mediante la evaluación conjunta de desempeño funcional y costos de implementación. Su cumplimiento corresponde a una selección experimental condicionada por la evidencia, no a la certificación de una alternativa que satisfaga todos los requisitos en cualquier escenario. El cierre de los objetivos específicos es el siguiente:

\begin{enumerate}
\item \textbf{Establecer modelos y conjuntos de referencia.} La metodología identifica las redes y separa los conjuntos de clasificación, validación estática y lazo cerrado. La identificación de las referencias permite reproducir las comparaciones dentro de cada banco. En radiografías, la diferencia de topología del PTQ INT8 y la reutilización del conjunto de 624 imágenes durante el entrenamiento limitan la atribución causal y la generalización externa.
\item \textbf{Generar y auditar las variantes.} Los resultados documentan Float32, PTQ INT8, QAT INT8, INT16x8 y precisión mixta, con la cobertura de la Tabla~\ref{tab:cobertura_variantes}. La generación y evaluación se acreditan para los modelos indicados; la auditoría de formas, tipos y parámetros de cuantización no es exhaustiva ni uniforme para todas las conversiones. Por ello, este objetivo tiene cumplimiento parcial en su componente de auditoría integral. La denominación INT16x8 no se equipara automáticamente a su ejecución nativa en destino.
\item \textbf{Evaluar los efectos numéricos y funcionales.} Las métricas de clasificación y las 200 trayectorias MPC permiten evaluar este objetivo en los bancos definidos. La precisión mixta recupera parte del desempeño perdido frente a INT8, pero en MPC no alcanza la referencia Float32. La concordancia numérica por inferencia y el éxito en lazo cerrado se mantienen como resultados distintos.
\item \textbf{Caracterizar ejecución, latencia y memoria.} Los perfiles y las validaciones describen las rutas CPU y NPU; los listados aportan evidencia estática complementaria de instrucciones. Se caracteriza la memoria del modelo y de la red compilada, sin extenderla a una medición uniforme del firmware completo. La comparación con L476 no aísla causalmente el aporte de Helium y no permite concluir ahorro energético.
\item \textbf{Implementar la validación UART y el banco HIL.} La demostración radiográfica y los ocho ensayos MPC acreditan la integración evaluada. En MPC, cinco ensayos cumplen el criterio terminal y tres no lo cumplen. Los ensayos acreditan el funcionamiento del banco con planta simulada, sin establecer robustez general, control de una planta física ni cumplimiento estricto de tiempo real.
\item \textbf{Comparar y justificar alternativas por aplicación.} La Tabla~\ref{tab:decision_final} fundamenta selecciones condicionadas. PTQ INT8/NPU se prioriza provisionalmente como compromiso radiográfico; QAT/NPU tiene la menor latencia medida y Mixed-5 conserva las decisiones Float32 en PC. En MPC, mixta/M55 es la alternativa parcialmente cuantizada investigada, mientras Float32/M55 conserva la referencia de desempeño. Esta distinción evita presentar la selección experimental como un óptimo universal.
\end{enumerate}

En conjunto, los hallazgos son compatibles con la hipótesis de que la sensibilidad de la tarea y la compatibilidad del grafo condicionan la selección más allá del número de bits. No constituyen una demostración causal aislada de cada factor. La brecha se aborda al contrastar la precisión solicitada con la ejecución obtenida y al vincular ambas con el desempeño de la aplicación. Permanecen abiertas la auditoría uniforme de todas las conversiones, la validación radiográfica externa y la comparación HIL de las alternativas bajo idénticas condiciones. Las conclusiones se circunscriben a las comparaciones experimentales documentadas.

